% Auto-generated from meeting_2567-06-07_stakeholder-survey-results.md (กบ.วช. format v2)
\documentclass{meetingminutes}
\meetingcommittee{คณะกรรมการบริหารหลักสูตร วิศวกรรมคอมพิวเตอร์และปัญญาประดิษฐ์ (บัณฑิตศึกษา)}
\meetingnumber{๓}{๒๕๖๗}
\meetingdate{วันที่ ๗ มิถุนายน พ.ศ. ๒๕๖๗}
\meetingplace{ห้องประชุมคณะเทคโนโลยีอุตสาหกรรม}
\meetingstart{๑๓.๐๐ น.}
\meetingend{๑๕.๑๐ น.}

\begin{document}
\maketitle

\psection{ผู้มาประชุม}
\begin{participants}
  \pmember{1}{รองศาสตราจารย์ ดร.อิสระ อินจันทร์}{ที่ปรึกษาหลักสูตร}{ประธาน}
  \pmember{2}{รองศาสตราจารย์ ดร.กันต์ อินทุวงศ์}{ที่ปรึกษาหลักสูตร}{รองประธาน}
  \pmember{3}{ผู้ช่วยศาสตราจารย์ ดร.พิศิษฐ์ นาคใจ}{อาจารย์ประจำหลักสูตร}{กรรมการ}
  \pmember{4}{ผู้ช่วยศาสตราจารย์ ดร.กาญจนา ดาวเด่น}{อาจารย์ประจำหลักสูตร (เลขานุการ)}{กรรมการ/เลขานุการ}
  \pmember{5}{ผู้ช่วยศาสตราจารย์ ดร.พัชรี มณีรัตน์}{อาจารย์ประจำหลักสูตร}{กรรมการ}
  \pmember{6}{ผู้ช่วยศาสตราจารย์ ดร.พิทักษ์ คล้ายชม}{รองคณบดีฝ่ายวิชาการ คณะเทคโนโลยีอุตสาหกรรม}{กรรมการ}
  \pmember{7}{ผู้ช่วยศาสตราจารย์ ดร.โสกณ วิริยะรัตนกุล}{อาจารย์ประจำหลักสูตร}{กรรมการ}
  \pmember{8}{ผู้ช่วยศาสตราจารย์ ดร.ธัชชัย อยู่ยิ่ง}{อาจารย์ประจำหลักสูตร}{กรรมการ}
  \pmember{9}{ผู้ช่วยศาสตราจารย์ ดร.อภิศักดิ์ พรหมฝ่าย}{อาจารย์ประจำหลักสูตร}{กรรมการ}
\end{participants}

\psection{ผู้ไม่มาประชุม}
\noindent ไม่มี\par

\startmeeting
\openingchair{รองศาสตราจารย์ ดร.อิสระ อินจันทร์}{กล่าวเปิดการประชุมและดำเนินการประชุมตามระเบียบวาระ ดังนี้}

\agenda{๑}{รับรองรายงานการประชุมครั้งที่ 2/2567}
ที่ประชุมรับรองรายงานการประชุมครั้งที่ 2/2567 โดยไม่มีการแก้ไข


\agenda{๒}{สรุปผลการสำรวจความต้องการผู้มีส่วนได้ส่วนเสีย (Stakeholder Needs Survey)}
ผศ.ดร.กาญจนา ดาวเด่น รายงานผลการสำรวจความต้องการผู้มีส่วนได้ส่วนเสีย ที่ดำเนินการระหว่างวันที่ 2–12 มิถุนายน 2567 ผ่าน Google Form:

สรุปผลสำคัญ (n = 31 คน):

\begin{longtable}{|p{7cm}|p{6cm}|}
  \hline
  \textbf{ประเด็น} & \textbf{ผล} \\
  \hline
  ต้องการศึกษาต่อระดับ ป.โท & 61.3\% (n=19) \\
  \hline
  ต้องการเรียนช่วง Weekend & 83.9\% (n=26) \\
  \hline
  เป้าหมายหลัก: ต่อยอดองค์ความรู้ & 45.2\% → PLO1, PLO3 \\
  \hline
  เป้าหมายรอง: ใช้ประกอบอาชีพ & 41.9\% → PLO4, PLO5 \\
  \hline
  เป้าหมาย: แก้ปัญหางาน/องค์กร & 22.6\% → PLO2 \\
  \hline
\end{longtable}

กลุ่มผู้มีส่วนได้ส่วนเสียภายนอกที่ตอบแบบสำรวจ: - การไฟฟ้าฝ่ายผลิตแห่งประเทศไทย (PEA) - Fujitsu Thailand - โรงพยาบาลอุตรดิตถ์ - มหาวิทยาลัยเทคโนโลยีราชมงคลล้านนา ตาก - มหาวิทยาลัยสุโขทัยธรรมาธิราช - ก.ศ.น. จังหวัดอุตรดิตถ์ - ราชธาณีวิทยาลัย, วิทยาลัยบึงพระ


\agenda{๓}{ตรวจสอบ Needs-to-PLO Mapping}
ที่ประชุมพิจารณาตาราง Needs-to-PLO Mapping:

\begin{longtable}{|p{5.5cm}|p{2.5cm}|p{4.5cm}|}
  \hline
  \textbf{ความต้องการของ SHs} & \textbf{ร้อยละ} & \textbf{PLO ที่สะท้อน} \\
  \hline
  ทักษะสร้างซอฟต์แวร์/AI/นวัตกรรม & 45.2\% & PLO1, PLO3 \\
  \hline
  ประยุกต์แก้ปัญหาองค์กร/ท้องถิ่น & 22.6\% & PLO2 \\
  \hline
  สร้างองค์ความรู้ใหม่/วิจัย & — & PLO3 \\
  \hline
  จริยธรรมวิชาชีพ & — & PLO4 \\
  \hline
  สื่อสาร ทีมงาน เรียนรู้ตลอดชีวิต & 41.9\% & PLO5 \\
  \hline
\end{longtable}

ที่ประชุมยืนยันว่า PLO ทั้ง 5 ข้อสะท้อนความต้องการของผู้มีส่วนได้ส่วนเสียครบถ้วน โดยเฉพาะ PLO1 และ PLO3 ที่ตอบสนองความต้องการหลัก (45.2\%)

\resolution{1. รับทราบผลการสำรวจและยืนยัน PLOs ที่มีอยู่สอดคล้องกับความต้องการของ SHs
2. มอบหมายให้จัดทำรายงานสรุปผลการสำรวจเป็นเอกสารหลักฐาน (AUNQA-1-4a/b)
3. นำข้อมูลนี้เป็น evidence สำหรับการรายงาน AUN-QA R1.4}

\agenda{๔}{ความก้าวหน้าการเตรียมเปิดรับนักศึกษา}
ประธานรายงานว่า สภามหาวิทยาลัยอนุมัติหลักสูตรแล้วในเดือนเมษายน 2568 และขณะนี้กำลังเตรียมการประชาสัมพันธ์เปิดรับนักศึกษาในปีการศึกษา 2568 ภาคเรียนที่ 2


\adjournmeeting

\begin{sigblock}
\typist{ผู้ช่วยศาสตราจารย์ ดร.กาญจนา ดาวเด่น}
\reviewer{รองศาสตราจารย์ ดร.อิสระ อินจันทร์}{กรรมการ}
\chairsig{รองศาสตราจารย์ ดร.อิสระ อินจันทร์}{ประธานการประชุม}
\end{sigblock}

\end{document}
